\subsection{Projective Temporal Ordering and One-Way Activation}
  \label{subsec:monotonicity-and-arrow-of-time}

  Temporal ordering in Cosmochrony is not postulated as a property of the
  substrate~$\chi$, nor derived from an intrinsic dynamics.
  It emerges as a structural consequence of the admissibility cascade.

  Along the breadth-first search cascade on the Cayley graphs of the Heisenberg
  group~$\mathrm{Heis}_3(\mathbb{Z}/q\mathbb{Z})$, define the \emph{cumulative
projected capacity}
  \begin{equation}
    I(n) \;=\; \sigma_{\mathrm{pair}}(0) - \sigma_{\mathrm{pair}}(n),
  \end{equation}
  where $\sigma_{\mathrm{pair}}(n) = \sigma_c(n)\cdot\sigma_{q-c}(n)$ is the
  residual projective capacity of the conjugate pair~$(c,q{-}c)$ at BFS
  depth~$n$~\cite{Beau2026a28}.
  The quantity~$I(n)$ measures the spectral structure stably projected up to
  depth~$n$.

  \begin{proposition}[Projective temporal ordering --- numerical
  form~\cite{Beau2026pto}]
  For $q \in \{29, 61, 101, 151\}$, $I(n)$ is strictly non-decreasing at every
  BFS step and for every conjugate pair~$(c,q{-}c)$, with no physically
  meaningful regression across~$169$ pairs and~$310$ post-burn-in steps.
  \end{proposition}

  This monotonicity defines an intrinsic ordering on effective states:
  \begin{equation}
    U_t \prec U_{t+1}
    \;\Longleftrightarrow\;
    I(U_t) \le I(U_{t+1}).
  \end{equation}
  The ordering is not a dynamics on~$\chi$ but a relational structure induced
  by the non-injective projection~$\Pi$: it acts entirely within the fixed
  observable rank established in~\cite{Beau2026a28}, without enlarging the
  observable degrees of freedom.

  \paragraph{One-way activation.}
    The observed monotonicity is underpinned by an empirical property of the
    cascade: spectral modes, once stably projected, do not deactivate.
    This \emph{one-way activation} property is a numerical result; its
    analytical derivation from the admissibility constraints of~$\Pi$ remains
    an open direction.

  \paragraph{Arrow of time and irreversibility.}
    Irreversibility follows directly from one-way activation: no admissible
    trajectory can decrease~$I(n)$.
    The arrow of time is identified with the directed growth of~$I(n)$, arising
    prior to any statistical or thermodynamic
    description~\cite{Prigogine1997,Rovelli1991}.
    The relation to thermodynamic irreversibility is discussed in
    Section~\ref{subsec:entropy-arrow-of-time}.

  \paragraph{Energy as residual projective capacity.}
    Energy is interpreted as the residual projective capacity of effective
    configurations: it measures the degree to which further accumulation of
    projected spectral structure remains admissible.
    Formally, $\sigma_{\mathrm{pair}}(n) = \sigma_{\mathrm{pair}}(0) - I(n)$
    is the remaining capacity at depth~$n$, decreasing monotonically as the
    cascade advances.
    Admissible ordering paths exclude any effective decrease of~$I(n)$, which
    would correspond to a restoration of projective capacity incompatible with
    the one-way activation property.

  \paragraph{Projectability and kinematic saturation.}
    \label{par:projectability-kinematics}
    A change of velocity corresponds to a modification of the relational
    coherence constraints maintained by the projection.
    As velocity increases, the informational demand on the projection grows,
    progressively saturating its capacity.
    The bound~$c_\chi$ is a structural limit beyond which no stable and
    globally consistent projection can be maintained.
    Approaching saturation, part of the relational content of~$\chi$ becomes
    inaccessible to the effective description, manifesting as time dilation,
    length contraction, and horizon formation.
    Relativistic kinematics thus emerges as a consequence of finite projection
    capacity.\footnote{%
  This kinematic saturation concerns the ordering and coherence capacity of
  the projection itself and should not be conflated with thermodynamic entropy
  production.}

  \paragraph{Planck scale and relativistic bounds as projection limits.}
    \label{par:planck-c-projective-limits}
    The constants~$c$ and~$h$ are interpreted as complementary manifestations of
    a finite resolution of the projection from~$\chi$ to effective observables.
    The bound~$c$ limits the maximal admissible rate at which relational ordering
    can be projected, while~$h$ sets a lower bound on the granularity with which
    the relational flux can be resolved.
    Relativistic and quantum constraints thus emerge as complementary facets of
    a single structural limitation: the finite capacity and resolution of the
    projection.
